\section{Closing an Epoch}
\label{sec:close}
An epoch closes in three steps: every outgoing validator freezes a report, the outgoing committee agrees on a checkpoint built from $N-f$ reports, and the next committee installs it. Figure~\ref{fig:close} follows Alice's account through one close.

\begin{figure}[t]
\centering
\begin{tikzpicture}[>=Stealth,font=\footnotesize,node distance=2mm and 3mm]
\matrix (m) [matrix of nodes,row sep=1.5mm,column sep=4mm,nodes={anchor=west,font=\footnotesize}]{
 \textbf{report $V_1$} & $A_6$: F & $A$: N & $B$: voted \\
 \textbf{report $V_2$} & $A_6$: F & $A$: N & $B$: voted \\
 \textbf{report $V_3$} & $A_6$: F & $A$: N & $B'$: voted \\
 \textbf{report $V_4$} & $A_6$: F & $A$: voted & $B$: voted \\
 \textbf{report $V_6$} & $A_6$: F & \textcolor{cBad}{$A'$: N} & \\
};
\node[below=3mm of m.south,anchor=north,align=left,text width=0.96\columnwidth,font=\footnotesize] (res) {%
\textbf{checkpoint $S_e$}: $A_6$ final (inherited); $A$ \emph{locked, notarized}
($V_1$, $V_2$, $V_3$ show it notarized; $V_6$'s claim for $A'$ has no four first votes behind it and is dropped);\\
$B$ \emph{locked for recovery, origin $e$} (three reporter first votes, no rival with three); $B'$ dropped.\\
Promotion: $A$ is the only block kept at position 7 and has an epoch-$e$ notarization $\Rightarrow$ $A$ final in $S_e$.};
\end{tikzpicture}
\caption{Five of six reports close epoch $e$ for Alice's account. F = final, N = notarized, voted = first-voted but not notarized (the report's second set).}
\label{fig:close}
\end{figure}

\para{Reports (\ref{r:report}).}
When a validator's local clock says the epoch is over, it stops first- and final-voting in it, and signs one header committing to two sets. The first is a tagged view of its ledger: every block it holds, tagged \emph{final} (by an inherited checkpoint or by a certificate it has seen), \emph{notarized} (four first votes seen, no finality), or \emph{inherited} (kept by the previous checkpoint under a recovery lock, nothing new seen). The second set lists every block the validator first-voted in this epoch that the first set does not already show as notarized or final. The two sets are frozen: later events change the validator's live view, never its report. The header carries only the hashes of the two sets; whoever needs the sets reconstructs them (Section~\ref{sec:model}). A report is \emph{usable} at a validator once both sets are reconstructed and every claim in them is backed by evidence that validator holds: four first votes behind every ``notarized'', a certificate behind every ``final'', the reporter's own signed first vote behind every entry of the second set. A Byzantine reporter can make its own report unusable, and nobody else's.

The second set has a subtle requirement, found by model checking. A reporter that first-voted a block the previous checkpoint already keeps under a recovery lock tags it \emph{inherited} in the first set. If the second set excluded every block of the first set, as our earlier definition did, this fresh vote would appear nowhere, and a latent fast certificate could vanish from the checkpoint. The second set therefore excludes only blocks tagged notarized or final.

\para{The recipe (\ref{r:build}).}
A candidate checkpoint names exactly $N-f$ usable reports from distinct members, in a fixed order, plus a manifest of the evidence behind their claims. From these inputs every validator computes the same state by the following recipe, which is deterministic and never consults the candidate's own certificate.
\begin{enumerate}
\item \emph{Keep finality.} Every block a selected report shows final, with its ancestors and the sends its receives depend on, is final.
\item \emph{Lock notarized blocks.} At each position, if a selected report shows a block notarized, backed by four first votes, keep it under a \emph{notarization lock}. Two conflicting blocks cannot both be notarized (Section~\ref{sec:epoch}), so there is at most one.
\item \emph{Lock for recovery.} If no block at the position is notarized, but exactly one block has first votes from $r$ (3) of the selected reporters, keep it under a \emph{recovery lock} with origin $e$. If several conflicting blocks reach $r$, which needs a reporter who voted both, nothing is locked and the position is voted afresh in the next epoch; this is safe because a block that can still become fast-final has no rival with $r$ supporters.
\item \emph{Carry old locks.} Every lock of the previous checkpoint is carried forward, unless a final block excludes its block, or, for a recovery lock of origin $o$, the manifest contains an exclusion witness of epoch $o$ for a conflicting block. Nothing else releases a lock; in particular a newer epoch number never does.
\item \emph{Promote.} For each account, starting at its last final block, walk up while the next block is the only one kept at its position and has an epoch-$e$ notarization of first votes; those blocks become final with the checkpoint. Promotion never resolves a fork and never promotes a block kept by a recovery lock alone.
\end{enumerate}
The thresholds make this safe (Theorem~\ref{thm:closure}): whatever the old epoch could still confirm is kept as final or under a lock, and a lock can only be released by evidence of its own origin epoch, which cannot exist if the block it protects can still be final.

\para{Why locks have an origin.}
A recovery lock says ``three validators first-voted this block in epoch $o$; a fast certificate with five epoch-$o$ votes may exist.'' Such a certificate is excluded exactly by four epoch-$o$ supporters for a rival, and by nothing else: four epoch-$(o+1)$ votes for the rival say nothing about epoch $o$. So the lock records $o$, and a block may carry locks of several origins, one per epoch in which it was voted but not resolved. Releasing a lock on the strength of a newer epoch's votes was our second wrong rule; it let two validators finalize different blocks.

\para{Eligibility, agreement, and installation (\ref{r:install}).}
A candidate is \emph{eligible} if its $N-f$ reports are usable, its manifest hash matches, and the recipe yields the state whose digest it names. The outgoing committee runs the agreement service with this predicate; the service may be any Byzantine agreement protocol with external validity. The decided value is small, $N-f$ report headers and two digests, whatever the number of accounts. The next committee verifies the decision proof, reconstructs the inputs, recomputes the state, and installs it: its final set is the union of its base and what it had already finalized, replay applies each operation once, and nothing is rolled back. There is no second vote; the successor only checks.

\para{Continuing from a kept block.}
A kept block gets an instance at its own position in the open epoch, and no rival is admitted there while a lock on it stands. Validators that voted the dropped rival in the closed epoch may vote the kept block now: the rule against switching (\ref{r:cross}) ends when the checkpoint is decided, because a decided checkpoint that kept $A$ proves that no closed-epoch certificate for a rival can form. So the kept block normally collects a full set of first votes and is final in the next epoch. If it does not, the next checkpoint keeps it again with a new origin, and its owner can extend it with a fresh child, whose confirmation confirms the kept block below it. A position at which nothing was kept is simply voted afresh. In our fork workloads nearly all split positions were resolved this way within an epoch or two; the few that stayed kept were kept identically on every validator (Section~\ref{sec:eval}).
